\section{TLB 与 \texttt{invlpg}}

\term{TLB}（Translation Lookaside Buffer）是 CPU 内部的\term{页表翻译缓存}。
每次虚拟地址翻译都需要两次内存访问（读 PDE + 读 PTE），TLB 将最近的翻译结果
缓存起来，命中时可以跳过页表查找。

\begin{figure}[H]
  \centering
  % figures/ch08/fig06-tlb-flow.tex — auto-extracted from chapter source
\begin{tikzpicture}[
  box/.style={draw, rounded corners, minimum width=2.5cm, minimum height=0.8cm, font=\small, align=center},
  arr/.style={-{Stealth[length=3pt]}, thick},
]
  \node[box, fill=blue!10]   (cpu)  at (0, 0)   {CPU 发出\\虚拟地址};
  \node[box, fill=green!15]  (tlb)  at (4, 0)   {TLB 查找};
  \node[box, fill=yellow!15] (hit)  at (8, 1)   {命中 → 物理地址};
  \node[box, fill=red!10]    (miss) at (8, -1)  {未命中 → 走页表};
  \node[box, fill=orange!10] (walk) at (12, -1) {PDE → PTE\\填入 TLB};

  \draw[arr] (cpu) -- (tlb);
  \draw[arr] (tlb) -- (hit)  node[midway, above, font=\scriptsize\itshape] {hit};
  \draw[arr] (tlb) -- (miss) node[midway, below, font=\scriptsize\itshape] {miss};
  \draw[arr] (miss) -- (walk);
  \draw[arr, dashed] (walk) |- (hit);
\end{tikzpicture}

  \caption{TLB 命中与未命中流程}
  \label{fig:tlb-flow}
\end{figure}

当我们修改了一个 PTE（改映射或取消映射），TLB 中可能还缓存着\textbf{过时}的翻译。
必须用 \texttt{invlpg} 指令使对应条目失效：

\codefile{src/kernel/arch/paging\_flush.asm}
\begin{lstlisting}[style=nasmstyle]
; void vmm_invlpg(uint32_t virt_addr);
vmm_invlpg:
    mov eax, [esp + 4]      ; eax = virt_addr
    invlpg [eax]            ; 使该虚拟地址的 TLB 条目失效
    ret
\end{lstlisting}

\begin{whybox}
为什么不每次都重新加载 \reg{CR3}（刷新\textbf{全部} TLB）？

重新加载 CR3 会使\textbf{所有}非 Global 的 TLB 条目失效，代价远大于
\texttt{invlpg} 只失效一条。在高频映射操作中（如缺页处理），
\texttt{invlpg} 的性能优势很明显。但在大批量拆除映射时
（如 \texttt{vmm\_unmap\_identity}），重新加载 CR3 反而更简单高效。
\end{whybox}

\begin{thinkbox}
\texttt{vmm\_unmap\_identity()} 清除了数百个 PDE 后调用
\texttt{vmm\_load\_page\_directory()} 重新加载 CR3。
如果改成对每个被清除的 PDE 覆盖的虚拟地址逐页 \texttt{invlpg}，
需要多少次 \texttt{invlpg} 调用？
（提示：每个 PDE 覆盖 1024 页。假设清除了 16 个 PDE\ldots）
\end{thinkbox}

% ============================================================================

